\documentclass[11pt]{article}
\usepackage[utf8]{inputenc}
\usepackage[T1]{fontenc}
\usepackage{mathptmx}
\usepackage[a4paper,margin=2.3cm]{geometry}
\usepackage{booktabs,tabularx,array}
\usepackage{enumitem}
\usepackage{xurl}
\usepackage[hidelinks]{hyperref}
\usepackage{caption}
\captionsetup{labelfont=bf,justification=raggedright,singlelinecheck=false,skip=6pt}
\renewcommand{\tablename}{Supplementary Table}
\renewcommand{\thetable}{S\arabic{table}}
\newcolumntype{L}{>{\raggedright\arraybackslash}X}
\setlength{\parindent}{0pt}
\setlength{\parskip}{6pt}
\renewcommand{\arraystretch}{1.18}

\begin{document}

{\Large\bfseries Supplementary material\par}
\vspace{4pt}
{\large Machine learning model for predicting the risk of childhood stunting in Peru\par}
\vspace{4pt}
{\small Abel Is-Boset Guevara-Huasco, Zulin Pamela Vallejos-Cotrina, Ver\'onica Lisseth Vergara-Rojas,
Ferdinand Edgardo Pineda-Ancco, Nemias Saboya-Rios\par}
\vspace{2pt}
\hrule
\vspace{6pt}

\textbf{Contents}
\begin{itemize}[nosep,leftmargin=1.2em]
  \item Supplementary Table S1. Predictor variables of the model grouped by domain.
  \item Supplementary Table S2. Configuration of the algorithms compared.
  \item Supplementary Table S3. Constructs and items of the relevance assessment instrument.
  \item Supplementary Text S1. Check of the variable selection.
\end{itemize}

%% ------------------------------------------------------------------ S1
\begin{table}[h]
\caption{Predictor variables of the model grouped by domain (45 variables; ENDES code in parentheses).}
\small
\begin{tabularx}{\linewidth}{@{}p{3.4cm}L@{}}
\toprule
Domain ($n$) & Variables \\
\midrule
Child (7) & Sex (HC27), Age in months (HC1), Birth weight (M19), Size at birth reported by the mother (M18),
Altitude-adjusted hemoglobin (HC56), Interval since the previous birth (HC63), Usual resident (HV102) \\
Mother, pregnancy and breastfeeding (11) & Mother's education (HC62), Prenatal visits (M14), Month of first prenatal
visit (M13), Days of iron during pregnancy (M46), Antiparasitic drugs during pregnancy (M60), Place of delivery (M15),
Cesarean delivery (M17), Postpartum vitamin A (M54), Hours until breastfeeding initiation (M34), Formula in the first
days (M55G), Duration of breastfeeding (M4) \\
Household and territory (8) & Wealth index (HV271), Area of residence (HV025), Altitude (HV040), Household members
(HV012), Children under five in the household (HV014), Age of household head (HV220), Sex of household head (HV219),
Traditional transport (SH77F) \\
Water and sanitation (10) & Water all day (SH42), Time to water source (HV204), Water treatment (HV237), Boiled water
(HV237A), Bottled water (HV237G), Water stored in container (SH48), Container type (SH49), Container with lid (SH50),
Shared toilet facility (HV225), Iodine in salt (HV234) \\
Housing (9) & Rooms (SH71), Bedrooms (HV216), Lighting source (SH70), Kitchen smoke outlet (HV240), Kitchen in separate
room (HV242), Other cooking fuel (SH63), Glass windows (SH76B), Wooden windows (SH76C), Windows with screens (SH76D) \\
\bottomrule
\end{tabularx}
\end{table}

%% ------------------------------------------------------------------ S2
\begin{table}[h]
\caption{Configuration of the algorithms compared. Hyperparameters were set in advance, without a systematic search;
those not listed keep the default values of each library. \texttt{random\_state = 42} in all models.}
\small
\begin{tabularx}{\linewidth}{@{}p{3.4cm}L@{}}
\toprule
Model & Configuration \\
\midrule
Logistic regression & L2 regularization; balanced class weights; maximum of 500 iterations \\
Decision tree (CART) & Maximum depth of 5; balanced class weights \\
XGBoost & 300 trees; \texttt{scale\_pos\_weight} equal to the ratio of negative to positive class;
\texttt{min\_child\_weight} $= 10^{-7}$ \\
CatBoost & 300 iterations; balanced class weights; native categorical variables \\
Neural network (MLP) & Three hidden layers of 200 neurons; up to 50 iterations with early stopping; no weighting \\
LightGBM & 300 trees; balanced class weights; native categorical variables \\
\bottomrule
\end{tabularx}

\vspace{3pt}
{\footnotesize Software: Python 3.12, LightGBM 4.6.0, XGBoost 3.2.0, CatBoost 1.2.10, scikit-learn 1.8.0 and SHAP 0.52.0.}
\end{table}

\clearpage
%% ------------------------------------------------------------------ S3
\begin{table}[h]
\caption{Constructs and items of the relevance assessment instrument. Items were rated on a five-point Likert scale
(1 = strongly disagree; 5 = strongly agree). The instrument was administered in Spanish; the items are translated
here.}
\small
\begin{tabularx}{\linewidth}{@{}p{3.2cm}p{1.65cm}L@{}}
\toprule
Construct & Source & Items \\
\midrule
Job relevance & TAM2 [1] &
\textbf{R1.} In my field of work, using a system like this would be important.\newline
\textbf{R2.} In my field of work, using a system like this would be relevant for deciding on nutritional interventions. \\
Output quality & TAM2 [1] &
\textbf{C1.} The risk factors it identifies by department are consistent with my professional experience.\newline
\textbf{C2.} I am confident that the risk levels assigned by the system are accurate. \\
Result demonstrability & TAM2 [1] &
\textbf{D1.} I understand why the system assigns a risk level to a specific child.\newline
\textbf{D2.} I could explain the results provided by the system to a colleague. \\
Perceived usefulness & TAM [2] &
\textbf{U1.} The system would allow me to prioritize which children to attend first when resources are limited.\newline
\textbf{U2.} Using it would save time compared with current methods of case identification.\newline
\textbf{U3.} Overall, it would improve the effectiveness of targeting work. \\
Perceived ease of use & TAM [2] &
\textbf{F1.} Entering a child's data into the system is easy.\newline
\textbf{F2.} The risk level and the departmental map are easy to read.\newline
\textbf{F3.} Learning to use the system would not require extensive training. \\
Perceptions of external control & TAM3 [3] &
\textbf{K1.} Except for hemoglobin, the data requested by the system are available in the records of my field of work.\newline
\textbf{K2.} Hemoglobin data are usually available for the children I attend.\newline
\textbf{K3.} My institution has the means (equipment, connectivity, staff) to use a tool like this. \\
Subjective norm & TAM2 [1] &
\textbf{S1.} My colleagues would consider the use of a tool like this appropriate.\newline
\textbf{S2.} The health authorities in my field would support its adoption. \\
Intention to use & TAM [2] &
\textbf{I1.} I would use this system if it were available in my field of work.\newline
\textbf{I2.} I would recommend its use to colleagues in nutrition or public health. \\
\midrule
Overall rating & --- & From 1 to 10, what potential do you assign to this tool for public health work? \\
Open-ended questions & --- &
(1) What additional variables or information should the system include?
(2) What limitations do you identify from your professional experience?
(3) In what specific scenario of your practice would you apply this tool?
(4) Recommendations for a next version. \\
\bottomrule
\end{tabularx}

\vspace{3pt}
{\footnotesize The form also recorded the evaluator's profession, academic degree, years of clinical experience,
institution and region; these data were used only to describe the panel.}
\end{table}

%% ------------------------------------------------------------------ Text S1
\section*{Supplementary Text S1. Check of the variable selection}

To rule out that the feature selection had removed useful information, the effect of adding back the excluded
variables was checked. Of the 78 columns available in the preprocessed dataset, 33 were not among the 45 selected.
These were ranked by their relationship with stunting: for categorical attributes, the range between the highest and
lowest prevalence among categories with at least 500 observations was used, whereas for continuous variables the
statistic $2\,|\mathrm{AUC} - 0.5|$ was used.

Four experimental configurations were compared: the 45-variable model and three that added the top-ranked variables
(water, sanitation and housing; territory; and both), keeping the partition, the weights, the deterministic seed and
the classifier hyperparameters constant.

The ranking identified ten attributes with a marked relationship with stunting, among them the predominant material of
the walls (HV214), floor (HV213) and roof (HV215); the source of drinking water (HV201); the department of residence
(HV024); the residual chlorine concentration in water (SH227); the cooking fuel (HV226); and the type of toilet facility
(HV205). These attributes make up the block of water, sanitation, hygiene and material housing conditions, recognized as
an underlying determinant of linear growth faltering [4,5] and found among the main predictors in studies based on DHS
data [6]. None of them constitutes information leakage, as they describe structural household conditions observable
before a home visit.

Adding the ten attributes did not improve any of the three configurations tested: the AUC-ROC (0.8295), sensitivity
(76.44\%) and precision (36.36\%) were identical to those of the 45-variable model, because the model did not use the
added variables. These attributes do contain information: a classifier trained only with the water, sanitation and
housing block achieved an AUC-ROC of 0.6895. Nor do they duplicate the wealth index, which explained between
$\eta^2 = 0.038$ and $\eta^2 = 0.098$ of the variance of each one. The 45-variable subset already represents water,
sanitation and housing conditions through nineteen correlated attributes (Supplementary Table S1), so the exclusion of
the candidate variables is a methodologically sound decision.

\section*{References}
\begin{enumerate}[label={[\arabic*]},leftmargin=2em,itemsep=2pt]\small
\item Venkatesh V, Davis FD. A theoretical extension of the technology acceptance model: four longitudinal field
studies. Manage Sci 2000;46:186--204. \url{https://doi.org/10.1287/mnsc.46.2.186.11926}.
\item Davis FD. Perceived usefulness, perceived ease of use, and user acceptance of information technology. MIS Q
1989;13:319--40. \url{https://doi.org/10.2307/249008}.
\item Venkatesh V, Bala H. Technology acceptance model 3 and a research agenda on interventions. Decis Sci
2008;39:273--315. \url{https://doi.org/10.1111/j.1540-5915.2008.00192.x}.
\item Black RE, Victora CG, Walker SP, Bhutta ZA, Christian P, de Onis M, et al. Maternal and child undernutrition and
overweight in low-income and middle-income countries. Lancet 2013;382:427--51.
\url{https://doi.org/10.1016/S0140-6736(13)60937-X}.
\item Victora CG, Christian P, Vidaletti LP, Gatica-Dom\'inguez G, Menon P, Black RE. Revisiting maternal and child
undernutrition in low-income and middle-income countries: variable progress towards an unfinished agenda. Lancet
2021;397:1388--99. \url{https://doi.org/10.1016/S0140-6736(21)00394-9}.
\item Bitew FH, Sparks CS, Nyarko SH. Machine learning algorithms for predicting undernutrition among under-five
children in Ethiopia. Public Health Nutr 2021;25:269--80. \url{https://doi.org/10.1017/S1368980021004262}.
\end{enumerate}

\end{document}
